%!TEX root = ../main.tex

\chapter{序論}
近年，地球温暖化対策としての温室効果ガス排出削減は世界的に課題となっている．日本政府は「2050年カーボンニュートラル」を掲げ，2025年2月にはエネルギー安定供給，経済成長，脱炭素の同時実現を目指す「GX2040ビジョン」を閣議決定した． 日本のエネルギー自給率は2023年度時点で15.3％とG7諸国で最低水準にあり，ロシアによるウクライナ侵略や中東情勢の緊迫化に伴う経済安全保障上のリスクが顕在化している．一次エネルギー供給の8割以上を海外からの輸入化石燃料に依存する現状において，国産の脱炭素エネルギー源である再生可能エネルギーの導入加速が急務となっている\cite{a1}．

再生可能エネルギーの中でも，太陽光発電は比較的リードタイムが短く，工場や公共施設等の屋根を活用した分散型電源として期待されている．しかし，太陽光発電は気象条件に依存して出力が変動するため，大量導入は電力系統の需給バランスを不安定化させる要因となる． このため，地域内での出力制御が増加しており，系統への負荷を抑えつつ再エネを有効活用するための「自家消費」や「需給近接型」の運用が重要視されている．特に，市場連動型の電力料金プランの普及に伴い，発電事業者自らが電力市場の需給に応じた供給計画を策定し，蓄電池等の調整力を活用して運用を高度化させることが求められている\cite{a1}．

電力グリッドの最適化に関しては，多様な数理的アプローチが提案されている． 岩瀬（2017）は，電気と熱の融通を考慮した分散型エネルギーグリッドの数理計画モデルを構築し，一次エネルギー消費量とCO2排出量の削減効果を定量的に評価した\cite{a3}．秋山（2013）は，東日本大震災後の電力供給の脆弱性を背景に，直流マイクログリッドにおける電力利用効率向上を目指し，据置型だけでなく移動型バッテリーを併用した混合整数計画モデルを提案した\cite{a4}．また，三浦（2015）は，兵庫県南あわじ市の沼島をモデルとした実証研究を補完するため，住民の電力消費行動を模擬するエージェントモデルとシステムモデルを統合したシミュレーションを行い，災害時の自立性やCO2削減効果を検証した\cite{a5}． これらの研究は，数理計画やシミュレーションを用いてシステムの有効性を示している．しかし，需要予測や市場単価が刻々と変動する実運用の現場において，限られた計算時間内でいかに迅速かつ確実に最適解を導出するかという，具体的な実装・適用フェーズにおける実用的な検討については，依然としてさらなる検討の余地がある．

本研究では，先行研究の知見を踏まえつつ，より実務的な運用手法の確立を目指す．具体的には，徳島県内の工場をモデルとした電力グリッドシステムを対象とし，市場連動型料金体系や自己託送制度を考慮した蓄電池運用の最適化モデルを構築する． 本研究の独自性は，線形計画法と動的計画法の2つの手法を同一モデルに適用し，その計算パフォーマンスや実用性を比較・検証する点にある．これにより，工場のエネルギーマネジメントにおける最適な数理的手法の選定指針を提示する．

本論文の構成は以下の通りである． まず第2章では，本研究の対象となる太陽光発電設備，蓄電池，および自己託送先を含む電力グリッドシステムの構成と，運用において満たすべき諸条件について述べる． 第3章では，連続値を扱うことで高精度な最適化が可能な線形計画モデルの定式化について，目的関数や制約条件の詳細を説明する． 第4章では，蓄電池の状態を離散化し，長期間の計画に対して効率的な計算が期待できる動的計画モデルの構築と，最適性の原理に基づくBellman方程式の適用について述べる． 第5章では，実データを用いた計算機実験の結果を示し，線形計画法と動的計画法の双方の観点から，計画期間や離散化数が計算時間に与える影響を比較・検討する． 
最後に第6章では，本研究で得られた知見を総括し，今後の課題と展望について述べる．